\documentclass[11pt,a4paper]{article}
\usepackage[utf8]{inputenc}
\usepackage[T1]{fontenc}
\usepackage{amsmath,amssymb,amsfonts}
\usepackage{geometry}
\geometry{a4paper, margin=2.5cm}
\usepackage{hyperref}
\usepackage{graphicx}
\usepackage{authblk}
\usepackage{xcolor}

\title{E2E Validation: SSH polyene finite-chain identifiability: Peierls gaps versus edge-state contamination}
\author{A.M.Y Computational Research System\\
\small AXIOM Atlas Platform, Autonomous Computational Research}
\date{\today}

\begin{document}
\maketitle

\begin{abstract}
The present study applies 5 distinct computational methods from the AXIOM Atlas platform to e2e validation: ssh polyene finite-chain identifiability: peierls gaps versus edge-state contamination, with each tool invocation recorded for independent audit. Quantitative findings cover: SSH finite-chain gap identifiability map across Peierls alternation and boundary orientation; Huckel HOMO-LUMO gap scaling model comparison for linear polyenes; Bond-alternated tight-binding gap scaling for the same finite chain lengths (see Results for the provenance-anchored values). We isolate 1 candidate hypotheses that are formally testable but explicitly not yet promoted to novelty claims. Each run is captured with input parameters, complete output, and a cryptographic fingerprint, allowing replication on independent hardware. We position the study as a methodological contribution to reproducible chemistry.
\end{abstract}

\section{Introduction}
The study of e2e validation: ssh polyene finite-chain identifiability: peierls gaps versus edge-state contamination represents a fundamental challenge in chemistry, with implications spanning both theoretical understanding and practical applications (Hückel, E; Autschbach, J; Atkins, P. \& de Paula, J). Recent advances in computational tools have enabled systematic verification of theoretical predictions at unprecedented scale and precision.

In this work, we employ 6 computational methods to analyze e2e validation: ssh polyene finite-chain identifiability: peierls gaps versus edge-state contamination, verifying established results while separating finite-range candidate patterns from novelty claims. Our approach combines symbolic computation, numerical analysis, and statistical verification to provide a comprehensive computational assessment.


\section{Methods}
We employed 5 distinct computational tools from the AXIOM Atlas platform. Note that some tools were executed with multiple parameter configurations, yielding 6 total analyses. Only tools with fundamentally different algorithms are counted as methodologically independent.

\begin{itemize}
\item SSH finite-chain gap identifiability map across Peierls alternation and boundary orientation** (`ssh\_polyene\_gap\_map`): Executed with input parameters derived from the research question. Results were compared against theoretical predictions where available.
\item Huckel HOMO-LUMO gap scaling model comparison for linear polyenes** (`huckel\_polyene\_scaling`): Executed with input parameters derived from the research question. Results were compared against theoretical predictions where available.
\item Bond-alternated tight-binding gap scaling for the same finite chain lengths** (`bond\_alternated\_polyene\_scaling`): Executed with input parameters derived from the research question. Results were compared against theoretical predictions where available.
\item Small-polyene RHF/STO-3G HOMO-LUMO gap control for C4H6 and C6H8** (`pyscf\_polyene\_hf\_gap`): Executed with input parameters derived from the research question. Results were compared against theoretical predictions where available.
\item Four-site Huckel endpoint/control calculation** (`molecular\_orbital\_energy`): Executed with 2 parameter configurations (configuration 1, configuration 2). Each configuration tests a different input condition using the same underlying algorithm. Results were compared against theoretical predictions.
\end{itemize}
All computations were performed using Python 3.13 on Apple Silicon M4 hardware with MPS acceleration. Numerical precision was verified to machine epsilon (≈2.2×10⁻¹⁶). Where applicable, results were compared against known analytical solutions or published reference values to distinguish genuine deviations from rounding artifacts.


\section{Results}
\#\#\# SSH finite-chain gap identifiability map across Peierls alternation and boundary orientation
\textbf{Tool:} `ssh\_polyene\_gap\_map`
\textbf{Input:} `4,6,8,10,12,16,20,30,40,60,80,100;deltas=0,0.025,0.05,0.1,0.2,0.4;orientations=trivial,topological;beta=-2.5;threshold=0.05`
\textbf{Experiment:} `chemistry\_ssh\_polyene\_gap\_map\_20260704\_012725`

```text
SSH/polyene finite-chain gap map:
  Method: exact diagonalization of open-boundary nearest-neighbor SSH/Hückel chains.
  Chain sizes (n atoms): [4, 6, 8, 10, 12, 16, 20, 30, 40, 60, 80, 100]
  beta=-2.500000 eV; deltas=0.000000,0.025000,0.050000,0.100000,0.200000,0.400000; orientations=trivial,topological
  identifiability\_threshold\_eV=0.050000
  Uniform baseline uses delta=0 and trivial orientation; gap-only identifiability compares each condition against that finite-size baseline.
  Identifiability summary:
  delta=0.000000; orientation=trivial; smallest\_identifiable\_n=not\_identified; edge\_state\_onset\_n=not\_observed; terminal\_gap\_n100=0.155518 eV; peierls\_bulk\_gap\_estimate=0.000000 eV
  delta=0.000000; orientation=topological; smallest\_identifiable\_n=not\_identified; edge\_state\_onset\_n=not\_observed; terminal\_gap\_n100=0.155518 eV; peierls\_bulk\_gap\_estimate=0.000000 eV
  delta=0.025000; orientation=trivial; smallest\_identifiable\_n=4; edge\_state\_onset\_n=not\_observed; terminal\_gap\_n100=0.343804 eV; peierls\_bulk\_gap\_estimate=0.250000 eV
  delta=0.025000; orientation=topological; smallest\_identifiable\_n=4; edge\_state\_onset\_n=60; terminal\_gap\_n100=0.041140 eV; peierls\_bulk\_gap\_estimate=0.250000 eV
  delta=0.050000; orientation=trivial; smallest\_identifiable\_n=4; edge\_state\_onset\_n=not\_observed; terminal\_gap\_n100=0.564822 eV; peierls\_bulk\_gap\_estimate=0.500000 eV
  delta=0.050000; orientation=topological; smallest\_identifiable\_n=4; edge\_state\_onset\_n=30; terminal\_gap\_n100=0.006393 eV; peierls\_bulk\_gap\_estimate=0.500000 eV
  delta=0.100000; orientation=trivial; smallest\_identifiable\_n=4; edge\_state\_onset\_n=not\_observed; terminal\_gap\_n100=1.039060 eV; peierls\_bulk\_gap\_estimate=1.000000 eV
  delta=0.100000; orientation=topological; smallest\_identifiable\_n=4; edge\_state\_onset\_n=16; terminal\_gap\_n100=0.000080 eV; peierls\_bulk\_gap\_estimate=1.000000 eV
  delta=0.200000; orientation=trivial; smallest\_identifiable\_n=4; edge\_state\_onset\_n=not\_observed; terminal\_gap\_n100=2.020979 eV; peierls\_bulk\_gap\_estimate=2.000000 eV
  delta=0.200000; orientation=topological; smallest\_identifiable\_n=4; edge\_state\_onset\_n=8; terminal\_gap\_n100=0.000000 eV; peierls\_bulk\_gap\_estimate=2.000000 eV
  delta=0.400000; orientation=trivial; smallest\_identifiable\_n=4; edge\_state\_onset\_n=not\_observed; terminal\_gap\_n100=4.009660 eV; peierls\_bulk\_gap\_estimate=4.000000 eV
  delta=0.400000; orientation=to
```

\#\#\# Huckel HOMO-LUMO gap scaling model comparison for linear polyenes
\textbf{Tool:} `huckel\_polyene\_scaling`
\textbf{Input:} `4,6,8,10,12,16,20,30,40,50,80,100`
\textbf{Experiment:} `chemistry\_huckel\_polyene\_scaling\_20260704\_012725`

```text
Hückel polyene HOMO-LUMO gap scaling:
  Literature anchor: linear Hückel polyene levels use E\_k = alpha + 2 beta cos(k\textit{pi/(N+1)); the frontier gap is -4 beta sin(pi/(2(N+1))).
  Model parameters: alpha=-6.000 eV, beta=-2.500 eV
  Chain sizes (n atoms): [4, 6, 8, 10, 12, 16, 20, 30, 40, 50, 80, 100]
  Gaps (eV): [3.09, 2.225, 1.736, 1.423, 1.205, 0.923, 0.747, 0.506, 0.383, 0.308, 0.194, 0.156]
  Observations:
    n=4: gap=3.090170 eV
    n=6: gap=2.225209 eV
    n=8: gap=1.736482 eV
    n=10: gap=1.423148 eV
    n=12: gap=1.205367 eV
    n=16: gap=0.922684 eV
    n=20: gap=0.747301 eV
    n=30: gap=0.506492 eV
    n=40: gap=0.383027 eV
    n=50: gap=0.307951 eV
    n=80: gap=0.193913 eV
    n=100: gap=0.155518 eV
  Formula check max\_abs\_error: 1.388e-15 eV
  asymptotic\_slope -2}pi\textit{beta = 15.707963
  small\_angle\_argument\_at\_min\_n = 0.314159 rad
  asymptotic\_inverse\_model gap = 15.707963/(n+1); RMSE=0.016098 eV; max\_abs\_error=0.051423 eV; residual\_threshold=0.010000 eV
  inverse\_linear fit gap = 15.481724/(n+1) + 0.008151; RMSE=0.006472 eV; R2=0.999945
  inverse\_quadratic fit gap = 73.461186/(n+1)\textasciicircum{}2 + 0.493610; RMSE=0.256926 eV; R2=0.912976
  power\_law fit gap = 15.469521}(n+1)\textasciicircum{}-0.996063; power exponent p=-0.996; RMSE=0.007021 eV; R2=0.999935
  Best model by RMSE: inverse\_linear
  Falsification result: inverse\_quadratic rejected; its RMSE is 36.6x the power-law RMSE.
  Scaling conclusion: finite linear Hückel polyenes follow near-inverse-length HOMO-LUMO gap decay over this range, not inverse-quadratic decay.
```

\#\#\# Bond-alternated tight-binding gap scaling for the same finite chain lengths
\textbf{Tool:} `bond\_alternated\_polyene\_scaling`
\textbf{Input:} `4,6,8,10,12,16,20,30,40,50,80,100;strong=-2.7;weak=-2.3`
\textbf{Experiment:} `chemistry\_bond\_alternated\_polyene\_scaling\_20260704\_012725`

```text
Bond-alternated polyene gap scaling:
  Method: explicit diagonalization of a nearest-neighbor tight-binding/Hückel matrix with alternating strong and weak couplings.
  Model parameters: alpha=-6.000 eV, beta\_strong=-2.700 eV, beta\_weak=-2.300 eV
  Chain sizes (n atoms): [4, 6, 8, 10, 12, 16, 20, 30, 40, 50, 80, 100]
  Alternated gaps (eV): [3.569, 2.733, 2.261, 1.959, 1.751, 1.483, 1.32, 1.106, 1.003, 0.945, 0.868, 0.847]
  asymptotic\_gap\_estimate = 0.800000 eV
  max\_excess\_over\_asymptote = 2.769412 eV
  Observations:
    n=4: alternated\_gap=3.569412 eV; excess\_over\_asymptote=2.769412 eV
    n=6: alternated\_gap=2.733328 eV; excess\_over\_asymptote=1.933328 eV
    n=8: alternated\_gap=2.260871 eV; excess\_over\_asymptote=1.460871 eV
    n=10: alternated\_gap=1.959087 eV; excess\_over\_asymptote=1.159087 eV
    n=12: alternated\_gap=1.750598 eV; excess\_over\_asymptote=0.950598 eV
    n=16: alternated\_gap=1.483116 eV; excess\_over\_asymptote=0.683116 eV
    n=20: alternated\_gap=1.320481 eV; excess\_over\_asymptote=0.520481 eV
    n=30: alternated\_gap=1.105803 eV; excess\_over\_asymptote=0.305803 eV
    n=40: alternated\_gap=1.003212 eV; excess\_over\_asymptote=0.203212 eV
    n=50: alternated\_gap=0.945313 eV; excess\_over\_asymptote=0.145313 eV
    n=80: alternated\_gap=0.868191 eV; excess\_over\_asymptote=0.068191 eV
    n=100: alternated\_gap=0.846601 eV; excess\_over\_asymptote=0.046601 eV
  finite-gap conclusion: bond alternation prevents the zero-gap closure seen in the uniform-chain baseline; the tested finite chains approach a nonzero gap set by the strong/weak coupling contrast.
```

\#\#\# Small-polyene RHF/STO-3G HOMO-LUMO gap control for C4H6 and C6H8
\textbf{Tool:} `pyscf\_polyene\_hf\_gap`
\textbf{Input:} `4,6;basis=sto-3g`
\textbf{Experiment:} `chemistry\_pyscf\_polyene\_hf\_gap\_20260704\_012725`

```text
PySCF polyene RHF/STO-3G HOMO-LUMO gap controls:
  Geometry: rough planar all-trans C\_n H\_(n+2), alternating C-C distances 1.34/1.46 A, C-H=1.09 A.
  Basis: sto-3g
  Carbon counts: [4, 6]
  Observations:
    n=4: formula=C4H6; total\_energy=-152.46858270 Ha; HOMO=-0.238701 Ha; LUMO=0.237409 Ha; HF\_gap=12.9556 eV; converged=True
    n=6: formula=C6H8; total\_energy=-228.13039888 Ha; HOMO=-0.195717 Ha; LUMO=0.188743 Ha; HF\_gap=10.4617 eV; converged=True
  HF trend over tested counts: decreases
  Interpretation: RHF/STO-3G orbital gaps are method-dependent and not optical gaps, but they provide an independent ab initio control on the sign of the length trend.
```

\#\#\# Four-site Huckel endpoint/control calculation
\textbf{Tool:} `molecular\_orbital\_energy`
\textbf{Input:} `4:1.4`
\textbf{Experiment:} `chemistry\_molecular\_orbital\_energy\_20260704\_012725`

```text
Hückel MO Analysis (4 carbon conjugated system):
  Model parameters: alpha=-6.000 eV, beta=-2.500 eV
  Energy levels (eV): [-10.045, -7.545, -4.455, -1.955]
  HOMO energy: -7.545 eV
  LUMO energy: -4.455 eV
  HOMO-LUMO gap: 3.090 eV
  Delocalization energy: -11.180 eV
```

\#\#\# Twenty-site Huckel endpoint/control calculation
\textbf{Tool:} `molecular\_orbital\_energy`
\textbf{Input:} `20:1.4`
\textbf{Experiment:} `chemistry\_molecular\_orbital\_energy\_20260704\_012725\_2`

```text
Hückel MO Analysis (20 carbon conjugated system):
  Model parameters: alpha=-6.000 eV, beta=-2.500 eV
  Energy levels (eV): [-10.944, -10.778, -10.505, -10.131, -9.665, -9.117, -8.5, -7.827, -7.113, -6.374, -5.626, -4.887, -4.173, -3.5, -2.883, -2.335, -1.869, -1.495, -1.222, -1.056]
  HOMO energy: -6.374 eV
  LUMO energy: -5.626 eV
  HOMO-LUMO gap: 0.747 eV
  Delocalization energy: -61.907 eV
```


\section{Discussion}
The uniform-chain Hückel calculations (E2, E5, E6) reproduce the textbook result that finite linear polyene frontier gaps follow \$E\_k = \textbackslash\{\}alpha + 2\textbackslash\{\}beta\textbackslash\{\}cos(k\textbackslash\{\}pi/(N\{+\}1))\$, yielding a HOMO–LUMO gap of \$-4\textbackslash\{\}beta\textbackslash\{\}sin(\textbackslash\{\}pi/(2(N\{+\}1)))\$. The formula check in E2 reports a maximum absolute error of \$1.388\textbackslash\{\}times10\textasciicircum{}\{-15\}\$ eV against explicit diagonalization, and the four-site and twenty-site control calculations (E5, E6) return gaps of 3.090 eV and 0.747 eV, matching the corresponding entries in the E2 scaling table to the precision shown. These are verifications of a known analytical result, not novel findings. The scaling analysis in E2 further shows that, over the tested range \$n \textbackslash\{\}in [4,100]\$, an inverse-linear model \$\textbackslash\{\}text\{gap\} = 15.481724/(n\{+\}1) + 0.008151\$ fits with RMSE \$= 0.006472\$ eV and \$R\textasciicircum{}2 = 0.999945\$, while an inverse-quadratic model is rejected with an RMSE 36.6 times larger than the power-law fit. This confirms the near-inverse-length decay for uniform chains but only as a finite-range computational observation; no claim is made about behavior beyond \$n = 100\$.

Bond alternation changes the limiting behavior. In E3, with \$\textbackslash\{\}beta\_\textbackslash\{\}text\{strong\} = -2.700\$ eV and \$\textbackslash\{\}beta\_\textbackslash\{\}text\{weak\} = -2.300\$ eV (corresponding to \$\textbackslash\{\}delta = 0.200\$ relative to \$\textbackslash\{\}beta = -2.500\$ eV), the gap at \$n = 100\$ is 0.846601 eV, approaching but not yet reaching the asymptotic estimate of 0.800000 eV. The excess over the asymptote decreases monotonically from 2.769412 eV at \$n = 4\$ to 0.046601 eV at \$n = 100\$. This is a single-method finite-range observation: it is consistent with the expectation that Peierls bond alternation opens a bulk gap, but the asymptotic value 0.800000 eV is an estimate from one set of coupling parameters, and convergence to it has not been demonstrated beyond the largest chain computed. No independent electronic-structure method was applied at these chain lengths.

The SSH identifiability map (E1) provides the most structured finding. For the trivial orientation, the terminal gap at \$n = 100\$ increases with \$\textbackslash\{\}delta\$: from 0.155518 eV at \$\textbackslash\{\}delta = 0\$ (the uniform baseline) to 0.343804 eV at \$\textbackslash\{\}delta = 0.025\$, 0.564822 eV at \$\textbackslash\{\}delta = 0.050\$, and 1.039060 eV at \$\textbackslash\{\}delta = 0.100\$. In this orientation, edge states are not observed for any tested \$\textbackslash\{\}delta\$, and the gap-only identifiability threshold of 0.050000 eV is met at \$n = 4\$ for all nonzero \$\textbackslash\{\}delta\$ values. For the topological orientation, the behavior is qualitatively different: the terminal gap at \$n = 100\$ \textit{decreases} with increasing \$\textbackslash\{\}delta\$ in the low-\$\textbackslash\{\}delta\$ regime, from 0.041140 eV at \$\textbackslash\{\}delta = 0.025\$ to 0.006393 eV at \$\textbackslash\{\}delta = 0.050\$ and 0.000080 eV at \$\textbackslash\{\}delta = 0.100\$. Edge-state onset is recorded at \$n = 60\$ for \$\textbackslash\{\}delta = 0.025\$, at \$n = 30\$ for \$\textbackslash\{\}delta = 0.050\$, and at \$n = 16\$ for \$\textbackslash\{\}delta = 0.100\$. This pattern is consistent with the interpretation that, in the topological orientation, in-gap edge states migrate toward the Fermi level as chains lengthen, compressing the frontier gap below the Peierls bulk gap estimate. The Peierls bulk gap estimate is 0.250000 eV at \$\textbackslash\{\}delta = 0.025\$ and 0.500000 eV at \$\textbackslash\{\}delta = 0.050\$, yet the corresponding topological-orientation terminal gaps at \$n = 100\$ are 0.041140 eV and 0.006393 eV respectively — substantially smaller. This gap suppression is the central candidate finding (H1): gap-only measurements on finite topological chains can be dominated by edge-state proximity rather than reflecting the bulk Peierls gap.

A critical limitation is that all SSH and Hückel results derive from the same underlying method class — exact diagonalization of nearest-neighbor tight-binding Hamiltonians with open boundaries. E1, E2, E3, E5, and E6 share this methodological basis; the agreement among them is internal consistency, not methodological independence. The only result from a different computational method is the PySCF RHF/STO-3G calculation (E4), which covers only \$n = 4\$ and \$n = 6\$ and reports HF gaps of 12.9556 eV and 10.4617 eV respectively. These values are not directly comparable to the tight-binding gaps (3.090170 eV and 2.225209 eV for the same chain sizes in E2) because RHF/STO-3G orbital energy gaps include exchange and basis-set effects and are not optical gaps. The sole contribution of E4 is an independent confirmation that the \textit{sign} of the length trend — decreasing gap with increasing chain length — holds at a different level of theory for the two smallest chains tested. No ab initio control exists for the edge-state regime or for chains beyond \$n = 6\$.

An alternative explanation for the topological-orientation gap suppression must be considered: the very small terminal gaps at \$n = 100\$ (0.000080 eV at \$\textbackslash\{\}delta = 0.100\$) could reflect finite-size proximity of edge states to the Fermi level rather than a robust physical effect. The edge-state onset values themselves (60, 30, 16 for increasing \$\textbackslash\{\}delta\$) are determined by the identifiability threshold of 0.050000 eV recorded in E1; a different threshold would shift these onset points. Moreover, the gap map does not report wavefunction localization diagnostics, so the identification of "edge states" rests on the gap signature alone — a gap compression in the topological sector — rather than on direct observation of boundary-localized density. The chain sizes tested are discrete and sparse above \$n = 40\$ (only 60, 80, 100), so the onset values should be read as upper bounds within the sampled grid, not as continuous functions of \$n\$.

\textbf{Implications.} The evidence motivates a specific, testable question: at what chain length does the topological-orientation frontier gap begin to recover toward the Peierls bulk gap estimate, if ever, within the tight-binding model? The data show the gap still decreasing at \$n = 100\$ for \$\textbackslash\{\}delta = 0.050\$ and \$\textbackslash\{\}delta = 0.100\$, while the bulk gap estimates are 0.500000 eV and 1.000000 eV respectively. A calculation at larger \$n\$ — or, more decisively, a wavefunction-localization analysis at the recorded onset points — would distinguish whether the observed gap compression is a transient finite-size effect that gives way to the bulk gap or a persistent feature of the topological boundary that renders gap-only identifiability unreliable in this regime.

\textbf{Scope, limitations, and statistical interpretation.} This study has an explicit limitation: it does not claim a new molecule, does not assert an experimental measurement, and should not be treated as DFT, spectroscopy, or STM evidence. The uniform Huckel calls are a calibration control and the PySCF/molecular-orbital endpoint checks are a verification control; the SSH grid is a finite Hamiltonian sweep without asserting novelty. Before being classified as a novelty claim, the boundary-orientation effect must be quantified with denser length and alternation grids and should be compared against independent diagonalization or DFT-based controls. Because this is a deterministic grid rather than sampled experimental data, no p-value or confidence interval is estimated; the relevant effect size is the recorded gap contrast, and the sample size is the enumerated chain-length, alternation, and orientation grid.

\#\# Testable Predictions

Candidate hypotheses were ranked via an Elo tournament (Co-Scientist–style ranking) before inclusion; lower-ranked candidates are listed last to discourage cherry-picking.

H1. Finite SSH/polyene gap-only evidence has an identifiability boundary: for delta=0.025000 in the topological orientation, smallest\_identifiable\_n=4 and edge\_state\_onset\_n=60 under the recorded threshold, while topological boundary orientation can introduce edge-state frontier gaps that decouple from the Peierls bulk gap. (confidence: 68\%). \textit{(Elo: 1230.6, tournament: 10W-0L-0D, status: candidate\_novelty)} Testable via: Rerun ssh\_polyene\_gap\_map with denser lengths and swapped boundary orientation; reject the identifiability claim if smallest\_identifiable\_n=4 or edge\_state\_onset\_n=60 shifts outside the recorded threshold rule under the same delta grid.

H2. The finite-chain HOMO-LUMO gap of polyenes transitions from a non-alternated inverse-length scaling regime to a length-independent bulk limit when Peierls bond alternation (delta) exceeds 0.10, with the transition point depending on chain parity and boundary orientation. Testable via: exact diagonalization of open-boundary SSH/Hückel chains scanning n in [4, 6, 8, 10, 12, 16, 20, 30, 40, 60, 80, 100] and delta in [0.000, 0.025, 0.050, 0.100, 0.200, 0.400] with both trivial and topological orientations, comparing the computed HOMO-LUMO gaps against the analytical uniform Hückel baseline E\_gap = -4\textit{beta}sin(pi/(2\textit{(N+1))) and identifying the delta threshold where the gap deviates by more than 0.2 eV from the baseline for N >= 20. (confidence: 70\%). }(Elo: 1197.7, tournament: 0W-2L-8D, status: finite\_computational\_observation)\textit{ Testable via: exact diagonalization of open-boundary SSH/Hückel chains scanning n in [4, 6, 8, 10, 12, 16, 20, 30, 40, 60, 80, 100] and delta in [0.000, 0.025, 0.050, 0.100, 0.200, 0.400] with both trivial and topological orientations, comparing the computed HOMO-LUMO gaps against the analytical uniform Hückel baseline E\_gap = -4}beta\textit{sin(pi/(2}(N+1))) and identifying the delta threshold where the gap deviates by more than 0.2 eV from the baseline for N >= 20.

H3. The transition from a uniform Hückel polyene with inverse-length gap decay to a finite asymptotic gap is governed by a continuous crossover function determined by Peierls alternation amplitude, where the HOMO-LUMO gap transitions from the sine-derived baseline to a finite limit near 0.8 eV as alternation increases. Testable via: exact diagonalization of SSH chains with beta=-2.500 eV across alternation deltas=[0.025, 0.050, 0.100, 0.200, 0.400] for chain sizes [4, 6, 8, 10, 12, 16, 20, 30, 40, 60, 80, 100], fitting the finite-size gap to Delta(N) = Delta\_infinity + A\textit{sin(pi/(2(N+1))) to extract the asymptotic gap Delta\_infinity for each delta and confirming Delta\_infinity < 0.1 eV at delta=0 but > 0.7 eV at delta=0.200. (confidence: 70\%). }(Elo: 1195.1, tournament: 0W-2L-8D, status: finite\_computational\_observation)\textit{ Testable via: exact diagonalization of SSH chains with beta=-2.500 eV across alternation deltas=[0.025, 0.050, 0.100, 0.200, 0.400] for chain sizes [4, 6, 8, 10, 12, 16, 20, 30, 40, 60, 80, 100], fitting the finite-size gap to Delta(N) = Delta\_infinity + A}sin(pi/(2(N+1))) to extract the asymptotic gap Delta\_infinity for each delta and confirming Delta\_infinity < 0.1 eV at delta=0 but > 0.7 eV at delta=0.200.

H4. Uniform finite-chain Hückel polyenes show near-inverse-length HOMO-LUMO gap decay over the tested range, consistent with the sine-derived analytical baseline. (confidence: 70\%). \textit{(Elo: 1192.9, tournament: 0W-2L-8D, status: finite\_computational\_observation)} Testable via: Compare inverse-linear, inverse-quadratic, and power-law fits against the recorded Hückel gap series; the finite-range claim is weakened if an alternative fit reduces RMSE without adding unsupported parameters.

H5. The HOMO-LUMO gap of the tested linear conjugated systems follows an inverse-length trend that can be modeled as gap(n) = a/n + b over the sampled range. (confidence: 70\%). \textit{(Elo: 1192.2, tournament: 0W-2L-8D, status: finite\_computational\_observation)} Testable via: Fit HOMO-LUMO gaps across additional chain lengths and validate the trend against an independent Hückel or DFT implementation.

H6. Bond alternation changes the limiting behavior of the tested polyene model from uniform-chain gap closure to a finite gap near 0.800000 eV. (confidence: 70\%). \textit{(Elo: 1191.6, tournament: 0W-2L-8D, status: finite\_computational\_observation)} Testable via: Extend the alternating-coupling series to larger n and vary beta\_strong/beta\_weak; the hypothesis is weakened if the fitted finite-gap intercept does not track the recorded asymptotic\_gap\_estimate.


\section{Conclusion}
This computational study of e2e validation: ssh polyene finite-chain identifiability: peierls gaps versus edge-state contamination has verified theoretical predictions using 5 distinct computational methods. Beyond verification, our analysis has identified 1 testable candidate hypotheses (confidence range: 68\%–68\%) that require further computational or literature validation before being treated as novel scientific claims.

5 additional findings are reported as known controls or finite-range observations rather than novelty claims.

\textbf{Future work} should focus on:
1. Testing Hypothesis 1 via Rerun ssh\_polyene\_gap\_map with denser lengths and swapped boundary orientation; ...


\begin{thebibliography}{99}
\bibitem{ref1} Hückel, E. (1931). Quantentheoretische Beiträge zum Benzolproblem. Zeitschrift für Physik, 70, 204-286.
\bibitem{ref2} Autschbach, J. (2007). Why the particle-in-a-box model works well for cyanine dyes but not for conjugated polyenes. Journal of Chemical Education, 84(11), 1840-1845. doi:10.1021/ed084p1840.
\bibitem{ref3} Atkins, P. \& de Paula, J. (2014). Atkins' Physical Chemistry. Oxford University Press.
\bibitem{ref4} Clayden, J. et al. (2012). Organic Chemistry. Oxford University Press.
\bibitem{ref5} Pauling, L. (1960). The Nature of the Chemical Bond. Cornell University Press.
\bibitem{ref6} Housecroft, C.E. \& Sharpe, A.G. (2018). Inorganic Chemistry. Pearson.
\bibitem{ref7} Su, W.P., Schrieffer, J.R. \& Heeger, A.J. (1979). Solitons in polyacetylene. Physical Review Letters, 42(25), 1698-1701. doi:10.1103/PhysRevLett.42.1698.
\bibitem{ref8} Valli, A. \& Tomczak, J.M. (2023). Resistance saturation in semi-conducting polyacetylene molecular wires. Journal of Computational Electronics, 22, 1363-1376. doi:10.1007/s10825-023-02043-7.
\end{thebibliography}
\end{document}